\documentclass[11pt]{letter}
\usepackage[utf8]{inputenc}
\usepackage[margin=2.5cm]{geometry}
\usepackage{hyperref}

\signature{Francisco Parra O., PhD\\Centro de Investigaci\'on en Tecnolog\'ias de Informaci\'on\\para la Agricultura y la Sociedad (CITIAPS)\\Santiago, Chile\\francisco.parra@citiaps.cl}

\address{}

\begin{document}

\begin{letter}{Editor-in-Chief\\Remote Sensing of Environment}

\opening{Dear Editor,}

I am submitting the manuscript entitled \textbf{``Symbolic Regression as Spectral Feature Engineering: Discovering Interpretable Indices for Hydrothermal Alteration Mapping from Sentinel-2''} for consideration as a Research Article in \textit{Remote Sensing of Environment}.

\textbf{The problem.} Spectral indices for hydrothermal alteration mapping were designed for ASTER over 20 years ago. With Sentinel-2 now the dominant sensor for geological remote sensing, there is a need for indices optimized to its spectral configuration---particularly its four Red Edge bands, which ASTER lacks.

\textbf{What we did.} We used symbolic regression (PySR) to automatically discover compact, interpretable spectral formulas from Sentinel-2 data across six hydrothermal alteration classes in Chile's Atacama Region, using SERNAGEOMIN field-mapped alteration polygons as ground truth.

\textbf{The key finding.} A Random Forest classifier trained on the six SR-derived features matches the performance of the same classifier using all 10 raw Sentinel-2 bands (mean AUC 0.898 vs.\ 0.899), demonstrating that symbolic regression effectively compresses 10 spectral bands into 6 interpretable features without measurable loss. Combining SR indices with classical spectral ratios yields the best overall performance (mean AUC = 0.920). This pattern replicates across four classifiers (RF, XGBoost, SVM-RBF, SVM-Linear) and two deep learning architectures (MLP, 1D-CNN).

\textbf{Cross-continent validation.} We validated the approach at the Cuprite mining district, Nevada (USA)---the most widely used benchmark for spectral alteration mapping---using USGS ASTER-derived ground truth. The feature engineering equivalence replicates (SR AUC = 0.942 vs.\ raw bands 0.940), confirming that the approach generalizes across geological provinces and continents.

\textbf{Why RSE.} This work bridges remote sensing methodology (automated spectral index discovery) with geological application (alteration mapping) and machine learning (interpretable feature engineering). The dual utility of SR indices---as standalone portable tools and as optimal features for ensemble classifiers---offers a new paradigm for sensor-specific spectral analysis that fits RSE's scope of advancing remote sensing science and its environmental applications.

\textbf{Novelty.} To our knowledge, this is the first application of symbolic regression to spectral index discovery for geological mapping. The demonstration that SR-derived features achieve near-lossless dimensionality reduction of Sentinel-2 bands is a general result applicable beyond hydrothermal alteration.

The manuscript has not been published elsewhere and is not under consideration by another journal. The author declares no competing interests.

\closing{Sincerely,}

\end{letter}
\end{document}
